\section{Results and Evidence Scope}
\label{sec:results}

This manuscript is a reproducible method and system report.  We do not fill the
joint-model tables with extrapolated or single-seed values.  At the repository
revision documented by this paper, the registered quantitative evidence
predates the combined global-query/reprojection experiment.  It provides
context for the backbone and dense tasks but cannot establish the benefit of
\cref{eq:total_loss}.

\paragraph{Revision snapshot.}
The exact source revisions used as evidence are:
\begin{center}
\small
\begin{tabular}{@{}ll@{}}
Scene/alignment/dataset & \texttt{179dac756aef137c9a35b1025ce76f0a31023648} \\
Camera-end canonicalization & \texttt{82fc58d5a677049e1d12d4867238af4ca22b9795} \\
Global query model & \texttt{2ae8cc3ea7ec6072d5595d691ff8c59a82f32ac7} \\
Query configuration routes & \texttt{077e84dc9d9edeeb4868ce8cff7c76e5d1ca8540} \\
\end{tabular}
\end{center}
The cross-branch algorithm is fully specified in this manuscript from
Issue~\#790, but this paper binds neither a frozen \#790 code revision nor a
numerical benchmark outcome.

\begin{table}[t]
  \centering
  \caption{Registered historical court-task evidence.  These runs predate and
  do not evaluate the combined global-query/reprojection method.}
  \label{tab:registered_evidence}
  \small
  \begin{tabularx}{\linewidth}{@{}p{0.27\linewidth}p{0.19\linewidth}X@{}}
    \toprule
    Registered run & Recorded value & Interpretation boundary \\
    \midrule
    KP14, 512-scale checkpoint & val. best 2.23 px; re-eval. 1.708886 px &
      Re-evaluation uses the validation source, not an independent test. \\
    Frozen \DINO{} segmentation & val. mIoU $\approx0.517$ &
      No tennis-domain self-supervised adaptation. \\
    Frozen adapted \DINO{} segmentation & val. mIoU $\approx0.800$ &
      Matched downstream setting; supports backbone adaptation only. \\
    B00 global subset alignment & 9.08 m / $89.94^\circ$ max drift &
      Negative run; rejected before dataset publication. \\
    B00 frozen holdout alignment & q95 1.240 m &
      Machine-rejected; a later user override is separate and remains marked. \\
    \bottomrule
  \end{tabularx}
\end{table}


The legacy \KP run reaches a registered validation best of 2.23 pixels and a
1.708886-pixel deterministic re-evaluation at 512-scale input.  The latter
loader reuses the validation source rather than an independent held-out test,
so we print it only as a software/data sanity check.  Likewise, the matched
frozen-backbone segmentation runs improve from approximately 0.517 to 0.800
validation mIoU after tennis-domain self-supervision.  These results support
the choice to expose \DINO{} features to the court model, but they neither
evaluate the pose query nor compare direct and joint objectives.
The corresponding repository records are:
\begin{quote}\small
\path{knowledge/nodes/run-i621-court-kp512-resume-r4.md}\\
\path{knowledge/nodes/run-i524-court-seg-baseline.md}\\
\path{knowledge/nodes/run-i524-court-seg-ssl.md}
\end{quote}
The KP record pins source commit
\begin{center}\texttt{41b318b39f757e09897964a32f28148a777b87f5}.\end{center}

Negative alignment evidence is equally informative.  In a historical B00
calibration attempt, aggregate fit residuals passed while subset rediscovery
shifted a court centre by up to 9.08 m and its orientation by 89.94 degrees.
The subsequent one-shot holdout achieved 0.832 weighted inliers but a 1.240 m
distance-weighted q95 and weak per-group coverage, so no machine-accepted
alignment/dataset contract was published.  A separate, explicit user-override
decision later published a SceneContract while preserving
\texttt{machine\_validation\_accepted=0}; that artifact is not evidence of a
passing automatic gate.  The three records are:
\begin{quote}\small
\path{knowledge/nodes/run-i618-b00-alignment-calibration-global-v1.md}\\
\path{knowledge/nodes/run-i618-b00-alignment-holdout-v1.md}\\
\path{knowledge/nodes/run-i618-b00-alignment-user-override-v1.md}
\end{quote}
This is the intended behavior of the fail-closed design: good
in-sample line support is not converted into thousands of confident but wrong
6DoF labels.

\paragraph{Claims reserved for registered runs.}
The blank numerical cells specified by \cref{tab:ablation_protocol} must be
populated from immutable run manifests containing commit, configuration,
checkpoint, split identities, seed, metrics, and prediction artifacts.  Until
those records exist, this paper makes no claim that joint-both improves a
baseline, that shortcut reliance is empirically eliminated, that a particular
decoder is Pareto optimal, or that the method is the production default.  The
method remains fully specified: a future run cannot change the oracle after
seeing its outcome.

\begin{table}[t]
  \centering
  \caption{\textbf{本稿の主張の境界。}
  実測した比較、実装済みの契約、
  これから検証する仮説、未測定の項目を分ける。
  }
  \label{tab:claim_scope}
  \small
  \begin{tabularx}{\linewidth}{@{}p{0.36\linewidth}p{0.14\linewidth}X@{}}
    \toprule
    主張 & 状態 & 根拠 \\
    \midrule
    失敗時に停止するメートル系アライメント & \statusimplemented &
      地上面推定、共通尺度の二面当てはめ、
      fit/holdout 分離、トポロジーと意味セグメントgate \\
    軌道生成とグループ分割 & \statusimplemented &
      凸包制約、弧長一様サンプル、trajectory group 分割 \\
    単一変換からの整合教師 & \statusimplemented &
      キーポイント・領域・線・線の意味・姿勢を同一変換から生成 \\
    検出器の並列姿勢回帰 & \statusimplemented &
      共有特徴からの密予測と姿勢10値の回帰 \\
    交差モデル比較（実写 validation / 合成 test） & \statusimplemented &
      domain ごと決定論的 120 件、共通 RANSAC 手続き、全推論 CPU \\
    視点拡張が実写汎化を改善する & \statushypothesis &
      同一枚数の狭域・広域比較と実写評価が必要 \\
    未知会場への汎化 & \statusunmeasured &
      会場を分けた評価集合が必要 \\
    3DGS 利用の因果効果 & \statusunmeasured &
      実写のみの同一アーキテクチャ比較が必要 \\
    採用視点分布の集計 & \statusunmeasured &
      学習へ渡った視点の方位・高さ・距離の集計図は未作成 \\
    \bottomrule
  \end{tabularx}
\end{table}

